\documentclass[11pt]{article}

\usepackage[a4paper,margin=25mm]{geometry}
\usepackage{amsmath,amssymb}
\usepackage{xcolor}
\usepackage{listings}

\definecolor{codeblue}{RGB}{25,70,120}
\definecolor{codegray}{RGB}{90,90,90}
\definecolor{codeback}{RGB}{247,248,250}

\lstdefinelanguage{pseudocode}{
  morekeywords={
    algorithm,input,output,require,if,then,else,end,for,each,while,do,
    return,true,false,unknown,and,or,not,in,with,continue
  },
  sensitive=false,
  morecomment=[l]{//},
  morestring=[b]"
}

\lstset{
  language=pseudocode,
  basicstyle=\ttfamily\small,
  keywordstyle=\color{codeblue}\bfseries,
  commentstyle=\color{codegray}\itshape,
  backgroundcolor=\color{codeback},
  frame=single,
  framesep=6pt,
  numbers=left,
  numberstyle=\tiny\color{codegray},
  numbersep=8pt,
  columns=fullflexible,
  keepspaces=true,
  breaklines=true,
  showstringspaces=false,
  escapeinside={(*@}{@*)}
}

\newcommand{\relu}{[\cdot]_+}
\newcommand{\R}{\mathbb{R}}
\newcommand{\True}{\mathsf{true}}
\newcommand{\False}{\mathsf{false}}
\newcommand{\Unknown}{\mathsf{unknown}}

\title{Pseudocode for Exact Symmetry and Identifiability Analysis\\
of Three-Layer Bottleneck ReLU Networks}
\author{Anonymous Authors}
\date{}

\begin{document}
\maketitle

\section{Scope and notation}

We follow the notation of the paper.  A three-layer architecture is
\[
  \mathcal A=(d,n_1,n_2,m),
\]
and its parameter space is denoted by \(\Theta_{\mathcal A}\).  A parameter is
\[
  \theta
  =
  \bigl(W^{(\ell)},b^{(\ell)}\bigr)_{\ell\in[3]}
  \in\Theta_{\mathcal A}.
\]
The two ReLU layers and the final affine layer are
\[
  \phi^{(\ell)}(x)
  =
  [W^{(\ell)}x+b^{(\ell)}]_+,\quad \ell\in[2],
  \qquad
  T^{(3)}(y)=W^{(3)}y+b^{(3)},
\]
and the realization is
\[
  f_\theta(x)
  =
  T^{(3)}\circ\phi^{(2)}\circ\phi^{(1)}(x).
\]
We use \(z^{(\ell,\theta)}\) and \(a^{(\ell,\theta)}\) for the preactivation
and activation maps, respectively, with \(a^{(0,\theta)}(x)=x\).

The fiber and generic fiber are
\[
  \mathcal S(\theta)
  :=
  \{\eta\in\Theta_{\mathcal A}:f_\eta=f_\theta\},
  \qquad
  \widetilde{\mathcal S}(\theta)
  :=
  \mathcal S(\theta)\cap\widetilde{\Theta}_{\mathcal A}.
\]
The relation \(\theta\sim\eta\) denotes equivalence under hidden-neuron
permutations and positive rescalings.  Thus \(\theta\) is identifiable when
\(f_\eta=f_\theta\) implies \(\eta\sim\theta\).

Throughout, \(\mathcal A=(d,n_1,n_2,m)\) is fixed with
\(d\geq n_1\geq2\), and all comparisons, signs, normalizations, and matrix
identities are exact.  The algorithms accept zero entries.  Every positive
equivalence or non-identifiability certificate is valid for arbitrary exact
parameters; completeness and identifiability are asserted only on
\(\widetilde{\Theta}_{\mathcal A}\).

Following Definition~7 of the paper, a second-layer neuron \(i\in[n_2]\)
\emph{hides hyperplane \(j\in[n_1]\)} if
\[
  W^{(2)}_{ij}>0,\qquad
  W^{(2)}_{ik}\leq 0\quad(k\neq j),\qquad
  b^{(2)}_i\leq 0.
\]
The parameter \(\theta\) hides hyperplane \(j\) if every \(i\in[n_2]\) does.
These weak inequalities include exact parameters with zeros.  On
\(\widetilde{\Theta}_{\mathcal A}\), the relevant signs are strict.

\section{Exact functional-equivalence certificate}

\begin{lstlisting}[caption={Exact equivalence dispatcher}]
algorithm EquivalenceTest((*@$\theta,\eta\in\Theta_{\mathcal A}$@*))
output: true if the certificate proves (*@$f_\theta=f_\eta$@*)

(*@$\bar\theta$@*) <- RemoveDeadSecondLayer((*@$\theta$@*))
(*@$\bar\eta$@*)   <- RemoveDeadSecondLayer((*@$\eta$@*))

if Width2((*@$\bar\theta$@*)) = Width2((*@$\bar\eta$@*)) = 0 then
    return (*@$\bar b^{(3,\theta)}=\bar b^{(3,\eta)}$@*)
else if exactly one reduced width is zero then
    return false
end if

(*@$j_\theta$@*) <- UniqueWeakHidingIndex((*@$\bar\theta$@*))
(*@$j_\eta$@*)   <- UniqueWeakHidingIndex((*@$\bar\eta$@*))

if exactly one of (*@$j_\theta,j_\eta$@*) is unknown then
    return false
end if

if both (*@$j_\theta,j_\eta$@*) are unknown then
    return NonHidingCertificate((*@$\bar\theta,\bar\eta$@*))
else if Width2((*@$\bar\theta$@*)) != Width2((*@$\bar\eta$@*)) then
    return false
else if Width2((*@$\bar\theta$@*)) >= 2 then
    return MultiHidingCertificate(
        (*@$\bar\theta,\bar\eta,j_\theta,j_\eta$@*))
else
    return WidthOneHidingCertificate(
        (*@$\bar\theta,\bar\eta,j_\theta,j_\eta$@*))
end if
\end{lstlisting}

For a non-hiding parameter, the induced one-hidden-layer map uses the paper's
notation
\[
  g_{\theta'}:\R_{\geq 0}^{n_1}\to\R^m,
  \qquad
  g_{\theta'}(y)
  =
  W^{(3)}[W^{(2)}y+b^{(2)}]_+ + b^{(3)},
\]
where \(\theta'=(W^{(2)},b^{(2)},W^{(3)},b^{(3)})\).

\begin{lstlisting}[caption={Non-hiding certificate}]
algorithm NonHidingCertificate((*@$\theta,\eta$@*))
((*@$\theta^\circ,\eta^\circ$@*)) <-
    CanonicalizeFirstLayerModuloPositiveScalingAndPermutation(
        (*@$\theta,\eta$@*))

if (*@$[W^{(1,\theta^\circ)},b^{(1,\theta^\circ)}]
       \neq [W^{(1,\eta^\circ)},b^{(1,\eta^\circ)}]$@*) then
    return false
end if

(*@$(\mathcal H_\theta,A_{R,\theta},b_{R,\theta},c_\theta)$@*) <-
    HyperplaneRepresentation((*@$g_{\theta'}$@*))
(*@$(\mathcal H_\eta,A_{R,\eta},b_{R,\eta},c_\eta)$@*) <-
    HyperplaneRepresentation((*@$g_{\eta'}$@*))

require (*@$\mathcal H_\theta=\mathcal H_\eta$@*)
require (*@$c_\theta(i)=c_\eta(i)$@*) on every (*@$H_i\in\mathcal H_\theta$@*)
require (*@$A_{R,\theta}=A_{R,\eta}$@*) and
        (*@$b_{R,\theta}=b_{R,\eta}$@*) on one common region (*@$R$@*)
return true exactly when all requirements hold
\end{lstlisting}

\begin{lstlisting}[caption={Hiding certificates}]
algorithm MultiHidingCertificate((*@$\theta,\eta,j_\theta,j_\eta$@*))
normalize the hidden first-layer direction
normalize and order all non-hidden affine hyperplanes
normalize every active second-layer row by its positive hidden coefficient
order second-layer rows canonically

require equal normalized first-layer directions
require equal non-hidden first-layer biases
require equal normalized (*@$W^{(2)},W^{(3)},b^{(3)}$@*)

(*@$t$@*) <- (*@$b^{(1,\eta)}_{j_\eta}-b^{(1,\theta)}_{j_\theta}$@*)
require (*@$b^{(2,\eta)}
          =b^{(2,\theta)}-tW^{(2,\theta)}_{:,j_\theta}$@*)
require (*@$t\geq
          \max_{i\in[n_2]} b^{(2,\theta)}_i/W^{(2,\theta)}_{ij_\theta}$@*)
return true exactly when all requirements hold

algorithm WidthOneHidingCertificate((*@$\theta,\eta,j_\theta,j_\eta$@*))
if MultiHidingCertificate((*@$\theta,\eta,j_\theta,j_\eta$@*)) = true then
    return true
end if

normalize the nonzero active row and output direction
match non-hidden first-layer hyperplanes up to unoriented scaling
recover the set of coordinates whose orientations are reversed
recover (*@$I\subseteq[n_1]\setminus\{j\}$@*) and apply (*@$M_I\in\mathcal G_j$@*)
require equality after the recovered sign flip, except for translation
recover (*@$t$@*) and require (*@$\eta\in T_j(M_I\cdot\theta)$@*)
return true exactly when all requirements hold
\end{lstlisting}

Here
\[
  T_j(\theta)
  =
  \left\{
    \left(
      W^{(1)},\
      b^{(1)}+te_j,\
      W^{(2)},\
      b^{(2)}-tW^{(2)}_{:,j},\
      W^{(3)},\
      b^{(3)}
    \right)
    :
    t\geq
    \max_{i\in[n_2]}
    \frac{b^{(2)}_i}{W^{(2)}_{ij}}
  \right\}
\]
is the paper's translation family, and for \(n_2=1\)
\[
  M_I
  =
  I_{n_1}
  -
  \sum_{k\in I}e_je_k^\top
  -
  2\sum_{k\in I}e_ke_k^\top,
  \qquad
  \mathcal G_j
  =
  \{M_I:I\subseteq[n_1]\setminus\{j\}\}.
\]

\section{Symmetry and identifiability classification}

\begin{lstlisting}[caption={Symmetry and identifiability classification}]
algorithm ClassifySymmetries((*@$\theta$@*))
input: (*@$\theta\in\Theta_{\mathcal A}$@*)
output: certified non-identifiable, generic-ID candidate,
        or undetermined

dead <- {(*@$i\in[n_2]$@*) :
         every entry of (*@$[W^{(2)}_{i,:},b^{(2)}_i]$@*) is nonpositive}
linear <- {(*@$i\in[n_2]$@*) :
           every entry of (*@$[W^{(2)}_{i,:},b^{(2)}_i]$@*) is nonnegative}
hiding <- {(*@$j\in[n_1]$@*) : (*@$\theta$@*) hides hyperplane (*@$j$@*)}

if dead (*@$\neq\varnothing$@*) or linear (*@$\neq\varnothing$@*)
   or hiding (*@$\neq\varnothing$@*) then
    return certified non-identifiable
else if SomeEntryIsZero((*@$[W^{(2)},b^{(2)}]$@*)) then
    return undetermined
else
    return generic-ID candidate
end if
\end{lstlisting}

On \(\widetilde{\Theta}_{\mathcal A}\), this is Theorem~13: non-identifiability
is equivalent to a strictly negative row, a strictly positive row, or a common
strict hiding column.  With zero entries, the weak conditions still give
explicit symmetries; otherwise the generic conclusion is left undetermined.

\section{Generic fiber dimension}

For the induced map \(g_{\theta'}:\R_{\geq0}^{n_1}\to\R^m\), let
\(\mathcal H=\{H_1,\ldots,H_k\}\) be its breakpoint-hyperplane arrangement and
let \(J=[n]\setminus[k]\) index the linear neurons, as in Proposition~5.  Thus
\(r:=|J|=n-k\).  The \emph{essential dimension} is \(\dim(V_S)\),
equivalently the dimension modulo \(\sim\).  The \emph{full dimension}
restores the \(n_1+n_2\) positive-scaling directions.

\begin{lstlisting}[caption={Dimension of the generic fiber modulo (*@$\sim$@*)}]
algorithm GenericFiberDimension((*@$\theta$@*))
input: generic (*@$\theta$@*); in the non-hiding case, one of Theorem 6's cases
((*@$\bar\theta,\mathrm{dead}$@*)) <- RemoveDeadSecondLayer((*@$\theta$@*))
(*@$q$@*) <- number of removed dead neurons

if Width2((*@$\bar\theta$@*)) = 0 then
    return (*@$\dim(\Theta_{\mathcal A})-m-(n_1+n_2)$@*)
end if

if (*@$\bar\theta$@*) hides hyperplane (*@$j$@*) then
    active_dimension <- (*@$\dim(T_j(\bar\theta)/{\sim})=1$@*)
else
    form the induced map (*@$g_{\bar\theta'}:\mathbb R_{\geq0}^{n_1}
                           \to\mathbb R^m$@*)
    (*@$k$@*) <- number of nonlinear neurons of (*@$g_{\bar\theta'}$@*)
    (*@$J$@*) <- linear-neuron set (*@$[n]\setminus[k]$@*)
    (*@$r$@*) <- (*@$|J|=n-k$@*)

    if (*@$r<\min(n_1,m)$@*) then
        active_dimension <- (*@$\dim(V_{S_{\bar\theta'}})=r^2$@*)
    else if (*@$r=n_1\leq m$@*) then
        active_dimension <- (*@$\dim(V_S)=r^2$@*)
    else if (*@$r=m<n_1$@*) then
        active_dimension <- (*@$\dim(V_S)=r^2$@*)
                            for every feasible orientation (*@$S$@*)
    else if (*@$r=m+1\leq n_1$@*) then
        active_dimension <- (*@$\dim(V_S)=m^2+m+n_1$@*)
    end if
end if

return active_dimension + (*@$q(n_1+m)$@*)
\end{lstlisting}

The cases for \(V_S\) are those of Theorem~6 applied to the induced shallow
map.  The full fiber dimension is obtained by restoring the \(n_1+n_2\)
positive-scaling directions.  This dimension statement is made under the
generic assumptions of the theorem.
ty
probability.

\end{document}
